\documentclass{situationist-preprint}

\title{Sufficiency Over Equality}
\msid{POL-08-a}
\msfamily{POL-08}
\msdiscipline{Political Philosophy}
\mstype{Principal manuscript}
\mskeywords{equality; sufficiency; levelling down; telic egalitarianism; priority view; difference principle; fitting attitudes}

\begin{document}
\maketitle

\begin{abstract}
Is equality valuable in itself? Telic egalitarians hold that it is: an outcome in which people are more equal is, to that extent, better. We argue that it is not. What matters is that everyone has enough, and the levelling-down objection is decisive against egalitarianism. The standard egalitarian reply concedes that levelling down is better \enquote{in one respect} and holds that it is worse all things considered. We show that the concession is the problem: on a fitting-attitude account of value, it implies that when misfortune makes the better off as badly off as the worse off, there is something in it to welcome, and there is not. The sufficiency view does not face the objection, since levelling down never brings anyone closer to having enough. We then make explicit how the result fits the modified difference principle of POL-01. That principle is not egalitarian, since it never ranks a levelling down above the status quo; it is prioritarian at the bottom; and it presupposes the sufficiency floor that our thesis identifies as what matters. The sufficiency thesis, in turn, is a thesis about what makes inequality objectionable, and it implies that inequality is not objectionable as such.
\end{abstract}

\section{The problem}

Many people think that inequality is bad: that a world in which some are much better off than others, through no fault or choice of their own, is worse for that reason. The thought can be developed in two ways \citep{Parfit1997}. \emph{Telic} egalitarians hold that inequality is bad in itself, so that an outcome is in one respect worse if it is more unequal; \citet{Temkin1993} gave the view its most careful defence. \emph{Deontic} egalitarians hold that we ought to promote equality for some other reason, typically because inequality results from injustice. Only the telic view holds that equality is valuable in itself.

\citet{Parfit1997} pressed against the telic view the \emph{levelling-down objection}. If inequality is bad in itself, then removing it by making the better off worse off, without making anyone better off, is in one respect an improvement. Suppose that some people are sighted and others blind; blinding the sighted would make everyone equal. Telic egalitarianism implies that this would be better in one respect, and many find that claim absurd. Parfit proposed instead the \emph{priority view}: benefiting people matters more the worse off they are. The priority view does not value equality, since it is concerned with each person's absolute level, not with relations between levels.

\citet{Frankfurt1987} had earlier argued that the moral appeal of equality is borrowed from a different ideal: what is important from the point of view of morality is not that everyone should have the same, but that each should have enough. The \emph{doctrine of sufficiency} was developed by \citet{Crisp2003}, who combined it with priority below the threshold, and criticized by \citet{Casal2007}, who distinguished its \emph{positive thesis}, that it is especially important that people live above a threshold, from its \emph{negative thesis}, that once everyone is above it no further distributive criteria apply. \citet{Shields2012} defended a version of sufficientarianism that accepts the positive thesis and rejects the negative one. \citet{Anderson1999} argued from a different direction that the point of equality is relational, a matter of people standing as equals, and that its distributive demands are sufficientarian.

We defend sufficiency over equality. Equality is not valuable in itself; what matters, in the place where egalitarians locate the value of equality, is that everyone has enough.

\section{Framework}

\begin{definition}\label{d:telic}
\emph{Telic egalitarianism} is the view that inequality between persons is bad in itself: of two outcomes, if one is more equal, it is better in one respect, whatever else is true of them.
\end{definition}

\begin{definition}\label{d:level}
A change is a \emph{levelling down} if it makes some people worse off, makes no one better off, and reduces inequality.
\end{definition}

\begin{definition}\label{d:threshold}
The \emph{sufficiency threshold} is the level of advantage adequate for a decent life and for the exercise of the basic liberties. The \emph{sufficiency thesis} has two parts: (i) it is of special and overriding importance that everyone be at or above the threshold; (ii) inequality is not objectionable as such, but only when, and because, it puts or keeps some below the threshold or undermines the conditions under which people stand as equals.
\end{definition}

Part (i) is Casal's positive thesis. Part (ii) is not her negative thesis: it does not say that nothing matters distributively above the threshold, but that what matters there is not equality. We return to this distinction in Section~4.

\begin{principle}[Fitting attitudes]\label{p:fitting}
An outcome is good in some respect only if some pro-attitude toward it in that respect, such as welcoming, being glad of or preferring it, is fitting.
\end{principle}

Principle~\ref{p:fitting} is the weak half of the fitting-attitude account of value. It requires only that value in a respect be the kind of thing to which a corresponding attitude could be fitting. It is accepted by most egalitarians, who hold that one should regret inequality and welcome its reduction.

\section{The argument}

\begin{lemma}\label{l:respect}
Telic egalitarianism implies that every levelling down is better in one respect.
\end{lemma}

\begin{proof}
A levelling down reduces inequality (Definition~\ref{d:level}). By Definition~\ref{d:telic}, the outcome with less inequality is better in one respect.
\end{proof}

Consider \emph{Natural Blinding}. In a population, some are sighted and others blind. An eruption leaves the sighted blind; no one is made better off, and no one is responsible. Natural Blinding is a levelling down, and by Lemma~\ref{l:respect} telic egalitarianism implies that it made things better in one respect.

\begin{theorem}\label{t:levelling}
Telic egalitarianism is false.
\end{theorem}

\begin{proof}
By Lemma~\ref{l:respect} and Principle~\ref{p:fitting}, if telic egalitarianism is true, some pro-attitude toward Natural Blinding is fitting: one should be, in some respect, glad of it. No such attitude is fitting. The formerly blind have gained nothing; they are not better placed, in their own lives or in their relations with others, than they were. A person who said that the eruption was terrible but that there was this to be said for it, that the sighted can no longer see what the blind cannot, would not be expressing a muted form of an appropriate attitude; she would be expressing an attitude that is unfitting, a form of spite. So telic egalitarianism is false.
\end{proof}

The case is natural so that no question of the agent's motive or of permissibility arises: the egalitarian cannot say that the act is condemned by a constraint, since there is no act. What remains is a claim about value, and it fails.

\begin{proposition}\label{p:suff}
The sufficiency thesis never favours a levelling down.
\end{proposition}

\begin{proof}
A levelling down makes no one better off, so it raises no one to the threshold. It may lower some below it, which part (i) of the thesis condemns. Reducing inequality is not, by part (ii), a respect in which an outcome is better.
\end{proof}

Theorem~\ref{t:levelling} and Proposition~\ref{p:suff} together establish that the objection is decisive against egalitarianism and does not touch sufficiency. It remains to explain the appeal of equality. The cases in which inequality seems clearly bad are cases in which some lack what others have in excess and lack enough: the hungry beside the rich. In these cases, part (i) of the sufficiency thesis condemns the situation, and it is the shortfall, not the gap, that does so; the same shortfall in a society where no one had more would be just as grave. Where inequality is objectionable among people who all have enough, it is because it confers power on some over others, or corrupts political equality, which part (ii) recognizes as derivative grounds \citep{Anderson1999}.

\section{Sufficiency and the difference principle}

POL-01 defends a modified difference principle: distributions are ranked by the level of the worst-off position, with ties broken by the total. It might be thought that a principle giving absolute priority to the worst off is egalitarian, and that the sufficiency thesis must reject it. Neither is the case.

\begin{proposition}\label{p:mdp}
The modified difference principle never ranks a levelling down above the status quo; it ranks every levelling down strictly below it.
\end{proposition}

\begin{proof}
In a levelling down no one is better off and someone is worse off, so the worst-off level does not rise and the total falls. If the worst-off level falls, the result is worse by the first clause of the principle. If it stays the same, the result is worse by the tie-breaker, since the total is lower.
\end{proof}

The modified principle is therefore not egalitarian in the sense of Definition~\ref{d:telic}: it assigns no value to the reduction of inequality and condemns every levelling down. It is prioritarian at the bottom, since it gives priority to the worst-off level, and its priority is grounded, in POL-01, in the reasoning of parties who choose under ignorance and cannot risk an inadequate life, not in the badness of the gap between the worst off and others.

The relation runs in both directions. First, the modified principle presupposes the sufficiency floor. The argument of POL-01 for choosing it relies on the assumption that some available principles guarantee every position a level adequate for a decent life and for the exercise of the basic liberties (Assumption~3.3 of POL-01-a); that level is the sufficiency threshold of Definition~\ref{d:threshold}. Where the assumption holds, maximizing the worst-off level secures that everyone has enough, so the principle satisfies part (i) of the sufficiency thesis. Second, the sufficiency thesis does not compete with the modified principle above the threshold. Part (ii) says that inequality is not objectionable as such; it does not say that nothing matters above the threshold. The modified principle permits any inequality that benefits the worst off, or that leaves them where they are and increases the total, which is what a principle should permit if inequality is not bad in itself. What the sufficiency thesis contributes is an account of why inequality, when it is objectionable, is objectionable: never because it is unequal. The sufficientarian complaint that the difference principle demands too much, mentioned in POL-01-a, comes from the negative thesis, which we reject; on the version defended here, sufficiency is the floor that the modified principle guarantees, and the principle's guidance above the floor is not in competition with it.

\section{Objections}

\begin{objection}[In one respect]
The egalitarian need not hold that levelling down is better, only that it is better in one respect and worse in others, and worse all things considered \citep{Temkin1993}. That claim is not absurd.
\end{objection}

\begin{reply}
The argument of Theorem~\ref{t:levelling} is directed at exactly this claim. It does not require that Natural Blinding be better all things considered, only that it be better in some respect, and it shows, by Principle~\ref{p:fitting}, that even this weak claim commits the egalitarian to a fitting attitude of gladness that is not fitting. The pluralist concession is what makes the view telic; without it there is no thesis that equality is valuable in itself.
\end{reply}

\begin{objection}[Indifference above the threshold]
Sufficientarianism is indifferent to inequalities above the threshold, however vast: between a world in which all have enough and a few have billions, and one in which all have a little more than enough \citep{Casal2007}.
\end{objection}

\begin{reply}
The objection applies to the negative thesis, which we reject. Above the threshold, the sufficiency thesis leaves room for other principles, including the modified difference principle, and it recognizes that great inequalities above the threshold are objectionable when they convert wealth into power over others or into unequal political influence. What it denies is that the vastness of the gap is in itself a reason against the first world.
\end{reply}

\begin{objection}[The threshold is arbitrary]
There is no non-arbitrary place to draw the line between enough and not enough, and a view that gives the line special importance inherits its arbitrariness.
\end{objection}

\begin{reply}
The threshold is fixed by its content, the level adequate for a decent life and the exercise of the basic liberties, and it is vague. Vagueness does not make it arbitrary. On the view of MET-18, vague predicates have sharp but unknowable boundaries; there is a level that is enough, though we cannot locate it exactly. In practice, policy can aim at a level safely above the region of uncertainty, as the floor of POL-01 does.
\end{reply}

\begin{objection}[Relational equality]
Equality matters in itself as a relation among persons: people should stand as equals. That value is not distributive, and the levelling-down objection does not touch it.
\end{objection}

\begin{reply}
We agree that relational standing matters, and part (ii) of the sufficiency thesis recognizes it. But its value lies in what it is for each person: freedom from domination, from contempt and from exclusion. Those are goods of persons, and each has a threshold, that of being able to stand as an equal \citep{Anderson1999}. The relational ideal is therefore an instance of the sufficiency thesis, not a case of equality valued in itself.
\end{reply}

\section{What the result does not establish}

The argument concerns telic egalitarianism. It does not refute deontic views on which inequality is objectionable because it results from injustice, which make no claim about the value of equality. Nor does it fix the threshold's level, or settle how benefits above it are to be distributed; the latter question is answered by the principles of POL-01.

\section{Conclusion}

Telic egalitarianism implies that a natural catastrophe which leaves the better off as badly off as the worse off has something to be said for it, and nothing can be said for it. Equality is therefore not valuable in itself. What matters is that everyone has enough, and inequality is objectionable only when it prevents this or undermines standing as equals. The modified difference principle of POL-01 agrees: it condemns every levelling down, gives priority to the worst off, and presupposes the floor that sufficiency requires.

\begin{thebibliography}{9}\small
\bibitem[Anderson(1999)]{Anderson1999} Anderson, E.~S. (1999). What is the point of equality? \emph{Ethics}, 109(2), 287--337.
\bibitem[Casal(2007)]{Casal2007} Casal, P. (2007). Why sufficiency is not enough. \emph{Ethics}, 117(2), 296--326.
\bibitem[Crisp(2003)]{Crisp2003} Crisp, R. (2003). Equality, priority, and compassion. \emph{Ethics}, 113(4), 745--763.
\bibitem[Frankfurt(1987)]{Frankfurt1987} Frankfurt, H. (1987). Equality as a moral ideal. \emph{Ethics}, 98(1), 21--43.
\bibitem[Parfit(1997)]{Parfit1997} Parfit, D. (1997). Equality and priority. \emph{Ratio}, 10(3), 202--221.
\bibitem[Shields(2012)]{Shields2012} Shields, L. (2012). The prospects for sufficientarianism. \emph{Utilitas}, 24(1), 101--117.
\bibitem[Temkin(1993)]{Temkin1993} Temkin, L.~S. (1993). \emph{Inequality}. Oxford University Press.
\end{thebibliography}

\end{document}
